\documentclass[12pt]{article}
\usepackage{fullpage,amsmath,amssymb,amsthm,hyperref,amsfonts}
\newtheorem{theorem}{Theorem}
\newtheorem{lemma}[theorem]{Lemma}
\newtheorem{proposition}[theorem]{Proposition}
\newtheorem{corollary}[theorem]{Corollary}
\theoremstyle{definition}
\newtheorem{definition}[theorem]{Definition}
\newtheorem{remark}[theorem]{Remark}
\newtheorem{construction}[theorem]{Construction}

\DeclareMathOperator{\rank}{rank}
\newcommand{\bR}{\mathbb{R}}

\title{Algebraic Relations on Determinantal Tensors\\
  (Solution to Problem 9)}
\author{}
\date{}

\begin{document}
\maketitle

\begin{abstract}
Let $n\geq 5$ and $A^{(1)},\ldots,A^{(n)}\in\bR^{3\times 4}$ be Zariski-generic.
Define $Q^{(\alpha\beta\gamma\delta)}\in\bR^{3\times 3\times 3\times 3}$ by
$Q^{(\alpha\beta\gamma\delta)}_{ijk\ell}=\det[A^{(\alpha)}_i;\,A^{(\beta)}_j;\,A^{(\gamma)}_k;\,A^{(\delta)}_\ell]$.
We prove that a polynomial map $\mathbf{F}:\bR^{81n^4}\to\bR^N$ satisfying (P1)--(P3) exists,
with all coordinate functions of degree exactly $5$ (independent of $n$).
\end{abstract}

\section{Setup and statement}

Write $[n]=\{1,\ldots,n\}$ and $[3]=\{1,2,3\}$. For $(\alpha,\beta,\gamma,\delta)\in[n]^4$ and
$I=(i,j,k,\ell)\in[3]^4$, the input is
\[
  P^{(\alpha\beta\gamma\delta)}_I
  \;:=\;
  \lambda_{\alpha\beta\gamma\delta}\,Q^{(\alpha\beta\gamma\delta)}_I,
\]
where $\lambda\in(\bR^*)^{[n]^4}$ (every entry nonzero by hypothesis) and $A^{(\alpha)}_i$
is the $i$-th row of $A^{(\alpha)}$.

Required properties of $\mathbf{F}:\bR^{81n^4}\to\bR^N$:
\begin{itemize}
\item[(P1)] $\mathbf{F}$ does not depend on $A^{(1)},\ldots,A^{(n)}$.
\item[(P2)] Degrees of coordinate functions are independent of $n$.
\item[(P3)] For Zariski-generic $A^{(\alpha)}$:
  $\mathbf{F}(\{P^{(\alpha\beta\gamma\delta)}_I\})=0$
  iff $\lambda_{\alpha\beta\gamma\delta}=u_\alpha v_\beta w_\gamma x_\delta$
  for some $u,v,w,x\in(\bR^*)^n$.
\end{itemize}

\begin{theorem}\label{thm:main}
Such $\mathbf{F}$ exists, with all coordinate functions of degree $5$.
\end{theorem}

\section{Construction}\label{sec:construction}

\begin{definition}
For each $s\in\{1,2,3,4\}$, the \emph{mode-$s$ unfolding matrix} $\mathbf{M}^{(s)}$ is the
$(3n)\times(3n)^3$ matrix whose rows are indexed by the pair $(\text{mode-}s\text{ global index}, \text{local index})$
and whose columns collect the remaining three mode pairs. Concretely:
\begin{equation}\label{eq:M1-def}
  \mathbf{M}^{(1)}_{(\alpha,i),\,(\beta,j,\gamma,k,\delta,\ell)} \;:=\; P^{(\alpha\beta\gamma\delta)}_{ijk\ell},
\end{equation}
\begin{equation}\label{eq:M2-def}
  \mathbf{M}^{(2)}_{(\beta,j),\,(\alpha,i,\gamma,k,\delta,\ell)} \;:=\; P^{(\alpha\beta\gamma\delta)}_{ijk\ell},
\end{equation}
\begin{equation}\label{eq:M3-def}
  \mathbf{M}^{(3)}_{(\gamma,k),\,(\alpha,i,\beta,j,\delta,\ell)} \;:=\; P^{(\alpha\beta\gamma\delta)}_{ijk\ell},
\end{equation}
\begin{equation}\label{eq:M4-def}
  \mathbf{M}^{(4)}_{(\delta,\ell),\,(\alpha,i,\beta,j,\gamma,k)} \;:=\; P^{(\alpha\beta\gamma\delta)}_{ijk\ell}.
\end{equation}
\end{definition}

\begin{construction}\label{const:F}
The map $\mathbf{F}$ consists of all $5\times5$ minors of $\mathbf{M}^{(s)}$ for $s=1,2,3,4$.
For $s=1$, every choice of $5$ row indices $\{(\alpha_a,i_a)\}_{a=1}^5$ and $5$ column indices
$\{(\beta_b,j_b,\gamma_b,k_b,\delta_b,\ell_b)\}_{b=1}^5$ contributes the minor
\begin{equation}\label{eq:minor-def}
  F^{(1)}_{\mathbf{r},\mathbf{c}}
  \;:=\;
  \det\!\bigl(P^{(\alpha_a\beta_b\gamma_b\delta_b)}_{i_a j_b k_b \ell_b}\bigr)_{1\leq a,b\leq 5}.
\end{equation}
Minors for modes $s=2,3,4$ are defined analogously by permuting the role of each mode,
using \eqref{eq:M2-def}--\eqref{eq:M4-def}.
\end{construction}

\section{Proof of Theorem~\ref{thm:main}}

\subsection{Properties (P1) and (P2)}

\textbf{(P1):} Each $F^{(s)}_{\mathbf{r},\mathbf{c}}$ in \eqref{eq:minor-def} is a polynomial
in the entries $\{P^{(\alpha\beta\gamma\delta)}_I\}$; no explicit reference to $A^{(\alpha)}$
appears.

\textbf{(P2):} Each coordinate is a $5\times5$ determinant, so has degree $5$ in the input,
independent of $n$.

\subsection{Mode-\texorpdfstring{$s$}{s} rank structure}

We analyse how the rank of each $\mathbf{M}^{(s)}$ depends on $\lambda$.

The cofactor expansion of $\det[A^{(\alpha)}_i;A^{(\beta)}_j;A^{(\gamma)}_k;A^{(\delta)}_\ell]$
along row $s$ is:
\begin{equation}\label{eq:cofactor-s}
  Q^{(\alpha\beta\gamma\delta)}_{ijk\ell}
  = \sum_{p=1}^{4}(-1)^{s+p}\,\bigl(A^{(X_s)}\bigr)_{\sigma_s p}\,
    M^{(s)}_p\!\bigl(A^{(X_1)}_{\sigma_1},A^{(X_2)}_{\sigma_2},A^{(X_3)}_{\sigma_3}\bigr),
\end{equation}
where $(X_1,\sigma_1,X_2,\sigma_2,X_3,\sigma_3,X_s,\sigma_s)$ refers to the remaining three rows
in the determinant after removing row $s$, and $M^{(s)}_p$ is the $3\times3$ minor obtained by
deleting column $p$.  Specifically, for mode $s=1$:
$(X_s,\sigma_s)=(\alpha,i)$ and the other three row-pairs are $(\beta,j),(\gamma,k),(\delta,\ell)$.

For each mode $s$ and each value of the mode-$s$ global index $\xi\in[n]$, define the
linear map $f^{(s,\xi)}:\bR^4\to\bR^{(3n)^3}$ by
\begin{equation}\label{eq:f-def-s}
  \bigl(f^{(s,\xi)}(c)\bigr)_{\text{column}}
  \;:=\;
  \lambda_{\xi^{(s)}}\,\det\!\bigl[c_s;\ldots\text{(rows from other modes)}\ldots\bigr],
\end{equation}
where the column index specifies the values of the other three mode pairs, $c_s=c\in\bR^4$
appears in row $s$, the remaining rows come from the corresponding $A^{(\cdot)}$ matrices,
and $\lambda_{\xi^{(s)}}$ denotes $\lambda$ with mode-$s$ index fixed to $\xi$ and others
from the column.  Then
\begin{equation}\label{eq:row-s}
  \bigl(\mathbf{M}^{(s)}\bigr)_{(\xi,\sigma),\,\cdot}
  \;=\; f^{(s,\xi)}\!\bigl(A^{(\xi)}_\sigma\bigr),
\end{equation}
because inserting $A^{(\xi)}_\sigma$ into row $s$ gives exactly the determinant entry of $P$.

The row space of $\mathbf{M}^{(s)}$ is $\sum_{\xi=1}^n f^{(s,\xi)}(V^{(s,\xi)})$, where
$V^{(s,\xi)}:=\mathrm{rowspace}(A^{(\xi)})\subseteq\bR^4$ is 3-dimensional for generic
$A^{(\xi)}$ with null vector $n^{(s,\xi)}\in\bR^4$.

\begin{lemma}[Mode-$s$ rank upper bound]\label{lem:rank-ub}
If $\lambda_{\alpha\beta\gamma\delta}=u_\alpha v_\beta w_\gamma x_\delta$, then
$\rank(\mathbf{M}^{(s)})\leq 4$ for every $s\in\{1,2,3,4\}$.
\end{lemma}

\begin{proof}
We give the proof for $s=1$; the other modes are identical by symmetry.
When $\lambda_{\alpha\beta\gamma\delta}=u_\alpha v_\beta w_\gamma x_\delta$:
\[
  P^{(\alpha\beta\gamma\delta)}_{ijk\ell}
  = u_\alpha\sum_{p=1}^4(-1)^{1+p}A^{(\alpha)}_{ip}
    \underbrace{v_\beta w_\gamma x_\delta M^{(1)}_p(A^{(\beta)}_j,A^{(\gamma)}_k,A^{(\delta)}_\ell)}_{=:(W_p)_{(\beta,j,\gamma,k,\delta,\ell)}}.
\]
The four vectors $W_p\in\bR^{(3n)^3}$ (for $p=1,2,3,4$) are independent of $\alpha$, and
every row $(\alpha,i)$ satisfies $\mathbf{M}^{(1)}_{(\alpha,i),\cdot}=u_\alpha\sum_p A^{(\alpha)}_{ip}W_p$.
Hence $\rank(\mathbf{M}^{(1)})\leq 4$.
\end{proof}

\begin{corollary}\label{cor:vanish}
$\lambda=u\otimes v\otimes w\otimes x\;\Longrightarrow\;\mathbf{F}=0$.
\end{corollary}

\begin{proof}
Each $\mathbf{M}^{(s)}$ has rank $\leq4$ by Lemma~\ref{lem:rank-ub}, so all its $5\times5$
minors vanish.
\end{proof}

\subsection{Rank lower bound and the converse}

The following single-mode lemma is the main technical input.

\begin{lemma}[Mode-$s$ rank lower bound]\label{lem:rank-lb-mode}
For Zariski-generic $A^{(1)},\ldots,A^{(n)}\in\bR^{3\times4}$ with $n\geq 5$, if
$\rank(\mathbf{M}^{(s)})\leq 4$, then $\lambda_{\alpha\beta\gamma\delta}$ is rank-1 in mode $s$:
there exist a scalar $\kappa_\xi\in\bR^*$ for each $\xi\in[n]$ and a function
$h:\,[n]^3\to\bR^*$ such that $\lambda_{\alpha\beta\gamma\delta}=\kappa_\xi \cdot h$ where
$(\xi,h)$ correspond to mode $s$.  Precisely:
\begin{align*}
  s=1: &\quad \lambda_{\alpha\beta\gamma\delta} = \kappa_\alpha\,h_{\beta\gamma\delta}, \\
  s=2: &\quad \lambda_{\alpha\beta\gamma\delta} = \kappa_\beta\,h_{\alpha\gamma\delta}, \\
  s=3: &\quad \lambda_{\alpha\beta\gamma\delta} = \kappa_\gamma\,h_{\alpha\beta\delta}, \\
  s=4: &\quad \lambda_{\alpha\beta\gamma\delta} = \kappa_\delta\,h_{\alpha\beta\gamma}.
\end{align*}
\end{lemma}

\begin{proof}
We prove the case $s=1$; the other cases follow by permuting roles of $\alpha,\beta,\gamma,\delta$.

Assume $\rank(\mathbf{M}^{(1)})\leq 4$; let $S\subseteq\bR^{(3n)^3}$ be the 4-dimensional row
space.  By \eqref{eq:row-s}:
\begin{equation}\label{eq:f-in-S}
  f^{(1,\alpha)}(V^{(1,\alpha)}) \;\subseteq\; S
  \qquad\text{for all }\alpha\in[n].
\end{equation}

\noindent\textbf{Step 1: A linear relation.}
Since $\dim S\leq4$, all the 3-dimensional subspaces $f^{(1,\alpha)}(V^{(1,\alpha)})\subseteq S$
lie in a common 4-dimensional space.

If the mode-1 unfolding of $\lambda$ (the $n\times n^3$ matrix $[\lambda_{\alpha\beta\gamma\delta}]_{\alpha,(\beta\gamma\delta)}$) already has rank~1, then $\lambda_{\alpha\beta\gamma\delta}=\kappa_\alpha h_{\beta\gamma\delta}$ and the conclusion \eqref{eq:mode1-rank1} holds.

Otherwise, the mode-1 unfolding has rank $\geq2$: pick $\alpha_1,\alpha_2\in[n]$ such that
$\lambda_{\alpha_1\beta\gamma\delta}\not\propto\lambda_{\alpha_2\beta\gamma\delta}$ as functions
of $(\beta,\gamma,\delta)$.  We claim $f^{(1,\alpha_1)}(V^{(1,\alpha_1)})\neq f^{(1,\alpha_2)}(V^{(1,\alpha_2)})$:
if they were equal, the coefficient-matching argument of Step~2 (applied to any
$c\in V^{(1,\alpha_1)}$ and the corresponding $c'\in V^{(1,\alpha_2)}$ with
$f^{(1,\alpha_1)}(c)=f^{(1,\alpha_2)}(c')$) would yield $\lambda_{\alpha_1\beta\gamma\delta}c
=\lambda_{\alpha_2\beta\gamma\delta}c'$ in $\bR^4$ for all $(\beta,\gamma,\delta)$;
since $c\neq0$ and $c'\neq0$, they are proportional ($c'=tc$ for a scalar $t$), giving
$\lambda_{\alpha_1\beta\gamma\delta}=t\lambda_{\alpha_2\beta\gamma\delta}$ for all
$(\beta,\gamma,\delta)$ — contradicting $\lambda_{\alpha_1}\not\propto\lambda_{\alpha_2}$.
Two distinct 3-dim subspaces of the 4-dim space $S$ span $S$:
\[
  f^{(1,\alpha_1)}(V^{(1,\alpha_1)}) + f^{(1,\alpha_2)}(V^{(1,\alpha_2)}) = S.
\]
For any $\alpha\in[n]$, the 3-dimensional image $f^{(1,\alpha)}(V^{(1,\alpha)})\subseteq S$
lies in this sum; hence for every $c\in V^{(1,\alpha)}$ there exist
$c_1,c_2\in\bR^4$ (depending linearly on $c$) with
\begin{equation}\label{eq:decomp1}
  f^{(1,\alpha)}(c) = f^{(1,\alpha_1)}(c_1) + f^{(1,\alpha_2)}(c_2).
\end{equation}

\noindent\textbf{Step 2: Coefficient matching.}
By the identity $\bigl(f^{(1,\xi)}(c)\bigr)_{(\beta,j,\gamma,k,\delta,\ell)}
= \lambda_{\xi\beta\gamma\delta}\det[c;A^{(\beta)}_j;A^{(\gamma)}_k;A^{(\delta)}_\ell]$
(a restatement of \eqref{eq:f-def-s} for $s=1$), equation~\eqref{eq:decomp1} gives:
\begin{equation}\label{eq:coeff-match1}
  \lambda_{\alpha\beta\gamma\delta}\det[c;\cdot\cdot\cdot]
  = \lambda_{\alpha_1\beta\gamma\delta}\det[c_1;\cdot\cdot\cdot]
  + \lambda_{\alpha_2\beta\gamma\delta}\det[c_2;\cdot\cdot\cdot]
\end{equation}
for each $(\beta,j,\gamma,k,\delta,\ell)$ (where $\cdot\cdot\cdot$ abbreviates
$A^{(\beta)}_j;A^{(\gamma)}_k;A^{(\delta)}_\ell$).
Fix $(\beta,\gamma,\delta)$ and let $(j,k,\ell)$ vary over $[3]^3$.
Write $\det[c';\cdot\cdot\cdot] = \langle c',\eta_{jk\ell}\rangle$ where
$\eta_{jk\ell}\in\bR^4$ is the vector of signed $3\times3$ cofactors from rows
$A^{(\beta)}_j,A^{(\gamma)}_k,A^{(\delta)}_\ell$ of the $4\times4$ matrix.
For Zariski-generic $A^{(\beta)},A^{(\gamma)},A^{(\delta)}$, the 27 vectors $\{\eta_{jk\ell}\}$
span $\bR^4$ (any $4\times4$ submatrix formed by picking four generic vectors from these
rows has full rank).
Equation~\eqref{eq:coeff-match1} thus becomes $\langle v,\eta_{jk\ell}\rangle=0$ for all
$(j,k,\ell)$, where $v:=\lambda_{\alpha\beta\gamma\delta}c-\lambda_{\alpha_1\beta\gamma\delta}c_1
-\lambda_{\alpha_2\beta\gamma\delta}c_2$.  Since $\{\eta_{jk\ell}\}$ spans $\bR^4$, this forces
\begin{equation}\label{eq:vec-eq1}
  \lambda_{\alpha\beta\gamma\delta}\,c
  = \lambda_{\alpha_1\beta\gamma\delta}\,c_1 + \lambda_{\alpha_2\beta\gamma\delta}\,c_2
  \quad\text{in }\bR^4,
\end{equation}
where $c_1,c_2$ are fixed (independent of $(\beta,\gamma,\delta)$) because the decomposition
\eqref{eq:decomp1} is a single vector identity in $\bR^{(3n)^3}$.

\noindent\textbf{Step 3: Proportionality.}
From \eqref{eq:vec-eq1} with pairs $(\beta_1,\gamma_1,\delta_1)$ and $(\beta_2,\gamma_2,\delta_2)$:
\begin{align}
  \lambda_{\alpha,1}\,c &= \lambda_{\alpha_1,1}\,c_1 + \lambda_{\alpha_2,1}\,c_2,
  \label{eq:A}\\
  \lambda_{\alpha,2}\,c &= \lambda_{\alpha_1,2}\,c_1 + \lambda_{\alpha_2,2}\,c_2,
  \label{eq:B}
\end{align}
with $\lambda_{\alpha,r}:=\lambda_{\alpha\beta_r\gamma_r\delta_r}$.
Multiply \eqref{eq:A} by $\lambda_{\alpha_2,2}$, \eqref{eq:B} by $\lambda_{\alpha_2,1}$, subtract:
\begin{equation}\label{eq:elim1}
  (\lambda_{\alpha,1}\lambda_{\alpha_2,2}-\lambda_{\alpha,2}\lambda_{\alpha_2,1})\,c
  = (\lambda_{\alpha_1,1}\lambda_{\alpha_2,2}-\lambda_{\alpha_1,2}\lambda_{\alpha_2,1})\,c_1.
\end{equation}

\noindent\textbf{Genericity rules out $c_1\parallel c$:}
If $c_1\parallel c$ for all $c\in V^{(1,\alpha)}$, then $V^{(1,\alpha)}\subseteq V^{(1,\alpha_1)}$;
since both have dimension $3$, this gives $V^{(1,\alpha)}=V^{(1,\alpha_1)}$, i.e.,
$n^{(1,\alpha)}\parallel n^{(1,\alpha_1)}$.  For Zariski-generic $A^{(\alpha)},A^{(\alpha_1)}$
this does not hold.  Hence for generic $A^{(\alpha)}$ and generic $c\in V^{(1,\alpha)}$,
$c_1\not\parallel c$.

\noindent\textbf{Consequence:} For any generic $c$ with $c_1\not\parallel c$, \eqref{eq:elim1}
and the linear independence of $c$ and $c_1$ force both scalar coefficients to vanish:
\begin{equation}\label{eq:2x2zero}
  \lambda_{\alpha,1}\lambda_{\alpha_2,2} = \lambda_{\alpha,2}\lambda_{\alpha_2,1}
  \quad\text{and}\quad
  \lambda_{\alpha_1,1}\lambda_{\alpha_2,2} = \lambda_{\alpha_1,2}\lambda_{\alpha_2,1}.
\end{equation}
The first equation gives
$\lambda_{\alpha\beta_1\gamma_1\delta_1}/\lambda_{\alpha\beta_2\gamma_2\delta_2}
=\lambda_{\alpha_2\beta_1\gamma_1\delta_1}/\lambda_{\alpha_2\beta_2\gamma_2\delta_2}$,
a ratio independent of $\alpha$.  Since this holds for every pair
$(\beta_1,\gamma_1,\delta_1),(\beta_2,\gamma_2,\delta_2)\in[n]^3$ (by repeating the
argument for each pair and noting that genericity of $A^{(\alpha)}$ ensures
$c_1\not\parallel c$ for that pair too), the $n^3\times n$ matrix
$[\lambda_{\alpha\beta\gamma\delta}]_{(\beta,\gamma,\delta),\alpha}$ has all $2\times2$ minors
equal to~$0$: it has rank~$1$.  Therefore
\begin{equation}\label{eq:mode1-rank1}
  \lambda_{\alpha\beta\gamma\delta} = \kappa_\alpha\,h_{\beta\gamma\delta}
\end{equation}
for some $\kappa_\alpha\in\bR^*$ and $h_{\beta\gamma\delta}\in\bR^*$.
\end{proof}

\begin{lemma}[Full rank-1]\label{lem:full-rank1}
If $\rank(\mathbf{M}^{(s)})\leq 4$ for all $s=1,2,3,4$, then $\lambda$ is rank-1.
\end{lemma}

\begin{proof}
Lemma~\ref{lem:rank-lb-mode} with $s=1$ gives $\lambda_{\alpha\beta\gamma\delta}=\kappa_\alpha h_{\beta\gamma\delta}$.
Applying Lemma~\ref{lem:rank-lb-mode} with $s=2$ to the tensor
$\lambda_{\alpha\beta\gamma\delta}=\kappa_\alpha h_{\beta\gamma\delta}$
(equivalently: the mode-2 rank condition forces the mode-2 unfolding to have rank~1):
$h_{\beta\gamma\delta}=\mu_\beta g_{\gamma\delta}$.
Hence $\lambda_{\alpha\beta\gamma\delta}=\kappa_\alpha\mu_\beta g_{\gamma\delta}$.
Applying mode~3: $g_{\gamma\delta}=\nu_\gamma x_\delta$.
Hence
\begin{equation}\label{eq:final-rank1}
  \lambda_{\alpha\beta\gamma\delta}
  = \underbrace{\kappa_\alpha}_{u_\alpha}
    \underbrace{\mu_\beta}_{v_\beta}
    \underbrace{\nu_\gamma}_{w_\gamma}
    \underbrace{x_\delta}_{x_\delta}.
\end{equation}
All factors are nonzero since $\lambda_{\alpha\beta\gamma\delta}\in\bR^*$.
\end{proof}

\subsection{Completing the proof of (P3)}

\begin{proof}[Proof of (P3)]
$(\Leftarrow)$: If $\lambda=u\otimes v\otimes w\otimes x$, then by Lemma~\ref{lem:rank-ub}
each $\mathbf{M}^{(s)}$ has rank $\leq4$, so all their $5\times5$ minors vanish, giving
$\mathbf{F}=0$.

$(\Rightarrow)$: If $\mathbf{F}=0$, then all $5\times5$ minors of each $\mathbf{M}^{(s)}$
vanish, i.e., $\rank(\mathbf{M}^{(s)})\leq4$ for $s=1,2,3,4$.  By Lemma~\ref{lem:full-rank1},
$\lambda$ is rank-1.
\end{proof}

Combining with the verified (P1) and (P2), Theorem~\ref{thm:main} is proved.\qed

\begin{remark}
The number of coordinate functions is $N = 4\binom{3n}{5}^2$
(four mode-unfoldings, each contributing $\binom{3n}{5}^2$ minors of size $5\times5$),
growing polynomially in $n$ while each coordinate has degree $5$, satisfying (P2).
\end{remark}

\begin{remark}
The condition $n\geq5$ ensures the null vectors $\{n^{(s,\xi)}\}_{\xi=1}^n$ of the generic
matrices $A^{(\xi)}\in\bR^{3\times4}$ are in general position in $\bR^4$ — specifically,
any two are non-proportional (holding on a Zariski-open set of the space of $n$-tuples of
$3\times4$ matrices, which is nonempty for $n\geq2$).  The spanning argument for
$\{\eta_{jk\ell}\}$ in Step~2 requires only generic $A^{(\beta)},A^{(\gamma)},A^{(\delta)}$
and holds for all $n\geq3$.
\end{remark}

\end{document}
